\documentclass[a4paper,12pt]{article}
\usepackage[utf8]{inputenc}
\usepackage[T1]{fontenc}
\usepackage[ngerman]{babel}
\usepackage{amsmath, amssymb}
\usepackage{geometry}
\usepackage{tikz}
\usetikzlibrary{shapes.geometric, arrows.meta, positioning}
\usepackage{parskip}
\usepackage{titlesec}
\usepackage{enumitem}

\geometry{left=3cm, right=3cm, top=2.5cm, bottom=2.5cm}
\titleformat{\section}{\large\bfseries}{}{0em}{}[\titlerule]
\titleformat{\subsection}{\normalsize\bfseries}{}{0em}{}

\tikzset{
  box/.style={rectangle, draw, rounded corners=3pt, text width=4.8cm,
              align=center, minimum height=1.0cm, font=\small, inner sep=4pt},
  arrow/.style={-{Latex[length=2mm]}, thick},
}

\begin{document}

\begin{center}
  {\LARGE\bfseries Studierhinweise zu Lektion 2}\\[0.5em]
  {\large Lineare Algebra (Modul 61112)}
\end{center}

\vspace{0.8em}

Diese Lektion liefert die konzeptuellen Grundlagen für den gesamten weiteren Verlauf des Moduls.
Nach einem Exkurs über das \textbf{Auswahlaxiom} und das \textbf{Lemma von Zorn} werden
bekannte Begriffe (Vektorräume, Basen, lineare Abbildungen) aus den Mathematischen Grundlagen
wiederholt und vertieft. Der Schwerpunkt liegt auf drei fundamentalen Konstruktionen:
\textbf{direkte Summe}, \textbf{Dualraum} und \textbf{Faktorraum}. Schließen Sie etwaige
Lücken bei Basiswechsel und Matrixdarstellungen frühzeitig, da diese Begriffe ständig benötigt werden.

% -----------------------------------------------------------------------
\section*{Struktur der Lektion 2}

\begin{center}
\begin{tikzpicture}[node distance=0.55cm]
  \node[box] (aus) {2.1 Auswahlaxiom\\Lemma von Zorn};
  \node[box, below=0.5cm of aus] (vr)  {2.2 Rückblick: VR, Basen,\\lineare Abbildungen,\\Matrixdarstellungen};
  \node[box, below=0.5cm of vr]  (ds)  {2.3 Direkte Summe\\Komplemente};
  \node[box, below=0.5cm of ds]  (du)  {2.4 Dualraum\\Dualbasen, duale Abbildungen};
  \node[box, below=0.5cm of du]  (fak) {2.5 Faktorräume*\\Konstruktion, Dimension};
  \node[box, below=0.5cm of fak] (äq)  {2.6 Äquivalenzrelationen\\Beispiel: äquivalente Matrizen};
  \draw[arrow] (aus) -- (vr);
  \draw[arrow] (vr) -- (ds);
  \draw[arrow] (ds) -- (du);
  \draw[arrow] (du) -- (fak);
  \draw[arrow] (fak) -- (äq);
\end{tikzpicture}
\end{center}

\newpage

% -----------------------------------------------------------------------
\section*{Zielelement 2.1 -- Auswahlaxiom und Lemma von Zorn}

\subsection*{Lerninhalte}
\begin{center}
\begin{tikzpicture}[node distance=0.55cm]
  \node[box] (ax)  {Auswahlaxiom (2.1.2)};
  \node[box, below=0.5cm of ax]  (po)  {Partielle Ordnungen,\\Ketten, obere Schranken};
  \node[box, below=0.5cm of po]  (zo)  {Lemma von Zorn (2.1.13)};
  \draw[arrow] (ax) -- (po);
  \draw[arrow] (po) -- (zo);
\end{tikzpicture}
\end{center}

\subsection*{Lernziele}
\begin{itemize}
  \item Das Auswahlaxiom und seine Bedeutung kennen (nicht klausurrelevant).
  \item Das Lemma von Zorn als Folgerung aus dem Auswahlaxiom einordnen.
  \item Verstehen, warum das Lemma von Zorn für den Basisergänzungssatz benötigt wird.
\end{itemize}

\subsection*{Selbstkontrollelement 2.1}
Was besagt das Lemma von Zorn informal? Welche Anwendung hat es im Basisergänzungssatz?

\newpage

% -----------------------------------------------------------------------
\section*{Zielelement 2.2 -- Rückblick: Vektorräume, Basen, lineare Abbildungen}

\subsection*{Lerninhalte}
\begin{center}
\begin{tikzpicture}[node distance=0.55cm]
  \node[box] (vr)  {Vektorraum $V$ über $K$,\\Unterraum, Basis};
  \node[box, below=0.5cm of vr]  (be)  {Basisergänzungssatz (2.2.11/12)\\auch für unendlich-dim. VR};
  \node[box, below=0.5cm of be]  (la)  {Lineare Abbildungen,\\Kern, Bild, Rang-Nullitätssatz};
  \node[box, below=0.5cm of la]  (md)  {Matrixdarstellungen (2.2.23)\\Basiswechsel / Transformationsformel};
  \draw[arrow] (vr) -- (be);
  \draw[arrow] (be) -- (la);
  \draw[arrow] (la) -- (md);
\end{tikzpicture}
\end{center}

\subsection*{Lernziele}
\begin{itemize}
  \item Vektorraum-Axiome, Unterraum-Kriterium und Dimensionsbegriff sicher beherrschen.
  \item Den Basisergänzungssatz kennen und anwenden.
  \item Matrixdarstellungen linearer Abbildungen berechnen; Transformationsformel (2.2.28) anwenden.
  \item Zusammenhang zwischen Matrixdarstellung und Basiswechsel verstehen.
\end{itemize}

\subsection*{Selbstkontrollelement 2.2}
Sei $f : \mathbb{R}^2 \to \mathbb{R}^2$, $f(x,y) = (x+y, 2x)$. Berechnen Sie die Matrixdarstellung von $f$ bzgl.\ der Standardbasis und der Basis $\{(1,1),(1,0)\}$.

\newpage

% -----------------------------------------------------------------------
\section*{Zielelement 2.3 -- Direkte Summe}

\subsection*{Lerninhalte}
\begin{center}
\begin{tikzpicture}[node distance=0.55cm]
  \node[box] (sum)  {Summe von Unterräumen $U_1 + \cdots + U_k$};
  \node[box, below=0.5cm of sum] (dir) {Direkte Summe $V = U_1 \oplus \cdots \oplus U_k$\\Charakterisierungen (Satz 2.3.6)};
  \node[box, below=0.5cm of dir] (dim) {Dimensionsformel (Satz 2.3.12)};
  \node[box, below=0.5cm of dim] (kom) {Komplement $W$: $V = U \oplus W$\\Existenz (Satz 2.3.17)};
  \draw[arrow] (sum) -- (dir);
  \draw[arrow] (dir) -- (dim);
  \draw[arrow] (dim) -- (kom);
\end{tikzpicture}
\end{center}

\subsection*{Lernziele}
\begin{itemize}
  \item Definition der direkten Summe kennen; Charakterisierungen aus Satz~2.3.6 anwenden.
  \item Dimensionsformel für direkte Summen benutzen.
  \item Den Begriff des Komplementes kennen; verstehen, dass zu jedem Unterraum ein Komplement existiert.
\end{itemize}

\subsection*{Selbstkontrollelement 2.3}
Sei $V = \mathbb{R}^3$ und $U = \langle e_1, e_2 \rangle$. Geben Sie ein Komplement von $U$ an und verifizieren Sie $V = U \oplus W$.

\newpage

% -----------------------------------------------------------------------
\section*{Zielelement 2.4 -- Dualraum}

\subsection*{Lerninhalte}
\begin{center}
\begin{tikzpicture}[node distance=0.55cm]
  \node[box] (lf)  {Linearformen $\varphi: V \to K$\\Dualraum $V^* = \mathrm{Hom}_K(V,K)$};
  \node[box, below=0.5cm of lf]  (db)  {Duale Basis $(v_1^*, \ldots, v_n^*)$\\zu einer Basis von $V$};
  \node[box, below=0.5cm of db]  (da)  {Duale Abbildung $f^* : W^* \to V^*$\\zu $f: V \to W$};
  \draw[arrow] (lf) -- (db);
  \draw[arrow] (db) -- (da);
\end{tikzpicture}
\end{center}

\subsection*{Lernziele}
\begin{itemize}
  \item Den Dualraum $V^*$ als Vektorraum aller Linearformen definieren und mit der Dualbasis rechnen.
  \item Koordinaten eines Vektors bzgl.\ der Dualbasis berechnen.
  \item Die duale Abbildung $f^*$ definieren und als lineare Abbildung erkennen.
\end{itemize}

\subsection*{Selbstkontrollelement 2.4}
Sei $V = \mathbb{R}^2$ mit Standardbasis $\{e_1, e_2\}$. Bestimmen Sie die duale Basis $\{e_1^*, e_2^*\}$ und berechnen Sie $e_1^*((3,7))$.

\newpage

% -----------------------------------------------------------------------
\section*{Zielelement 2.5/2.6 -- Faktorräume und Äquivalenzrelationen}

\subsection*{Lerninhalte}
\begin{center}
\begin{tikzpicture}[node distance=0.55cm]
  \node[box] (äq)  {Äquivalenzrelation auf $M$\\Reflexivität, Symmetrie, Transitivität};
  \node[box, below=0.5cm of äq]  (äkl) {Äquivalenzklassen $[m]$\\Quotientenmenge $M/{\sim}$};
  \node[box, below=0.5cm of äkl] (fak) {Faktorraum $V/U$\\Nebenklassen $v + U$};
  \node[box, below=0.5cm of fak] (dim) {$\dim(V/U) = \dim V - \dim U$};
  \draw[arrow] (äq) -- (äkl);
  \draw[arrow] (äkl) -- (fak);
  \draw[arrow] (fak) -- (dim);
\end{tikzpicture}
\end{center}

\subsection*{Lernziele}
\begin{itemize}
  \item Eine Äquivalenzrelation definieren und Äquivalenzklassen bestimmen.
  \item Den Faktorraum $V/U$ als Quotient eines Vektorraums nach einem Unterraum konstruieren.
  \item Die Dimension von $V/U$ berechnen.
  \item Das Beispiel der äquivalenten Matrizen (2.6) kennen.
\end{itemize}

\subsection*{Selbstkontrollelement 2.5/2.6}
Sei $V = \mathbb{R}^3$, $U = \langle e_3 \rangle$. Beschreiben Sie die Nebenklassen im Faktorraum $V/U$ geometrisch. Was ist $\dim(V/U)$?

\end{document}
